\documentclass[12pt]{article}
\usepackage{amsmath, amssymb}
\usepackage{geometry}
\usepackage{enumitem}
\usepackage{verbatim}
\geometry{margin=1in}

\title{CSE400 -- Fundamentals of Probability in Computing\\
Lecture 10: Randomized Min-Cut Algorithm}
\author{Instructor: Dhaval Patel, PhD}
\date{February 5, 2026}

\begin{document}

\maketitle

\hrule
\vspace{0.5cm}

\section{Min-Cut Problem}

\subsection{Why Use Min-Cut?}

The lecture states that the min-cut algorithm is used in applications related to:

\begin{itemize}
    \item Network connectivity
    \item Reliability
    \item Optimization
\end{itemize}

Specifically:

\begin{enumerate}
    \item \textbf{Network Design}
    \begin{itemize}
        \item Min-cut helps in improving the efficiency of communication.
        \item It is used to optimize network flow.
        \item The algorithm is used to find the minimum capacity cut.
    \end{itemize}

    \item \textbf{Communication Networks}
    \begin{itemize}
        \item It helps in understanding the vulnerability of networks to failures.
        \item It supports building robust and fault-tolerant communication networks.
    \end{itemize}

    \item \textbf{VLSI Design}
    \begin{itemize}
        \item In Very Large Scale Integration (VLSI) design, min-cut is useful for partitioning circuits into smaller components.
        \item This reduces interconnectivity complexity.
    \end{itemize}
\end{enumerate}

(Reference: Section 1.5, \textit{Application: A Randomized Min-Cut Algorithm in Probability and Computing (2nd Edition)}, as mentioned in the slide.)

\subsection{What is Min-Cut?}

\subsubsection*{Definition: Cut-Set}

A \textbf{cut-set} in a graph is:

\begin{quote}
A set of edges whose removal breaks the graph into two or more connected components.
\end{quote}

\subsubsection*{Definition: Min-Cut Problem}

Given a graph
\[
G = (V, E)
\]
with $n$ vertices,

The \textbf{minimum cut (min-cut) problem} is:

\begin{quote}
To find a minimum cardinality cut-set in $G$.
\end{quote}

Thus:

\begin{itemize}
    \item Input: Graph $G = (V, E)$
    \item Output: A cut-set with minimum number of edges
\end{itemize}

\subsection{Randomness in Min-Cut Algorithms}

The slides state:

\begin{itemize}
    \item Min-cut algorithms such as Karger’s algorithm are random.
    \item They can be sensitive to the initial choice of edges.
    \item If the algorithm contracts critical edges early, it might find a smaller cut.
\end{itemize}

Thus, correctness depends on contraction choices.

\subsection{Edge Contraction (Main Operation)}

The main operation in the min-cut algorithm is:

\begin{quote}
\textbf{Edge contraction}
\end{quote}

\subsubsection*{Definition: Edge Contraction}

Edge contraction is:

\begin{quote}
An operation that removes an edge from a graph while simultaneously merging the two vertices.
\end{quote}

\subsubsection*{Detailed Description (from slide)}

When contracting an edge $(u, v)$:

\begin{enumerate}
    \item Merge vertices $u$ and $v$ into a single vertex.
    \item Eliminate all edges connecting $u$ and $v$.
    \item Retain all other edges in the graph.
    \item The new graph:
    \begin{itemize}
        \item May have parallel edges.
        \item Has no self-loops.
    \end{itemize}
\end{enumerate}

This operation is illustrated in the diagram on pages 15--16 of the slides.

\section{Successful and Unsuccessful Min-Cut Runs}

\subsection{Successful Min-Cut Run}

Definition:

\begin{quote}
A successful min-cut run refers to the success in the outcome of an algorithm designed to find the minimum cut in a graph.
\end{quote}

The figure labeled “Successful Run” illustrates a sequence of contractions that preserve the minimum cut until the final two vertices remain.

Thus:

\begin{itemize}
    \item The algorithm contracts edges.
    \item The true minimum cut is not destroyed.
    \item The final cut corresponds to the minimum cut.
\end{itemize}

\subsection{Unsuccessful Min-Cut Run}

Definition:

\begin{quote}
An unsuccessful min-cut run refers to an iteration of a min-cut algorithm where the algorithm fails to correctly identify the minimum cut of a given graph.
\end{quote}

The slide explicitly states:

\begin{quote}
The algorithm fails to correctly identify the minimum cut.
\end{quote}

The figure labeled “Unsuccessful Run” shows contraction of edges that eliminate the true minimum cut.

Thus:

\begin{itemize}
    \item A critical edge is contracted.
    \item The minimum cut is destroyed.
    \item The final result is not the true minimum cut.
\end{itemize}

\section{Max-Flow Min-Cut Theorem}

The theorem states:

\begin{quote}
In a flow network, the maximum amount of flow passing from the source to the sink is equal to the total weight of the edges in a minimum cut.
\end{quote}

\subsection*{Definitions from Slide}

\begin{enumerate}
    \item \textbf{Capacity of a cut} \\
    = The sum of the capacity of edges in the cut that are oriented from a vertex $\in X$ to a vertex $\in Y$.

    \item \textbf{Minimum cut} \\
    = The cut in the network that has the smallest possible capacity.

    \item \textbf{Minimum cut capacity} \\
    = The capacity of the minimum cut.

    \item \textbf{Maximum flow} \\
    = The largest possible flow from source $(S)$ to sink $(T)$.
\end{enumerate}

Thus:
\[
\text{Maximum Flow} = \text{Minimum Cut Capacity}
\]

\section{Deterministic Min-Cut Algorithm}

\subsection{Stoer--Wagner Min-Cut Algorithm}

Let $s$ and $t$ be two vertices of graph $G$. \\
Let $G/\{s,t\}$ be the graph obtained by merging $s$ and $t$.

The theorem states:

\begin{quote}
A minimum cut of $G$ can be obtained by taking the smaller of:
\begin{itemize}
    \item A minimum $s$-$t$-cut of $G$, and
    \item A minimum cut of $G/\{s,t\}$.
\end{itemize}
\end{quote}

Reason (as given in slide):

Either:

\begin{enumerate}
    \item There exists a minimum cut that separates $s$ and $t$. \\
    $\rightarrow$ Then a minimum $s$-$t$-cut of $G$ is a minimum cut of $G$.

    \item There is no such cut separating $s$ and $t$. \\
    $\rightarrow$ Then a minimum cut of $G/\{s,t\}$ gives the answer.
\end{enumerate}

\subsection{Pseudocode (As Given)}

\subsubsection*{Algorithm 1: MinimumCutPhase(G, a)}

\begin{verbatim}
A ← {a};
while A ≠ V do
    add to A the most tightly connected vertex;
return the cut weight as the "cut of the phase";
\end{verbatim}

\subsubsection*{Algorithm 2: MinimumCut(G)}

\begin{verbatim}
while |V| ≥ 1 do
    choose any a from V;
    MinimumCutPhase(G, a);
    if the cut-of-the-phase is lighter than the current minimum cut then
        store the cut-of-the-phase as the current minimum cut;
    shrink G by merging the two vertices added last;
return the minimum cut;
\end{verbatim}

\section{Randomized Min-Cut Algorithm}

\subsection{Why Randomized Algorithm?}

The slides state:

Randomized algorithms:

\begin{itemize}
    \item Provide a probabilistic guarantee of success.
    \item Provide a more accurate estimate of the minimum cut with fewer iterations.
\end{itemize}

Additional properties listed:

\begin{itemize}
    \item Efficiency
    \item Parallelization
    \item Approximation Guarantees
    \item Avoidance of Worst-Case Instances
    \item Heuristic Nature
    \item Robustness
\end{itemize}

\subsection{Karger’s Randomized Algorithm}

The diagram shows:

\begin{itemize}
    \item Random selection of an edge.
    \item Edge contraction.
    \item Repetition until two vertices remain.
    \item The remaining parallel edges form the cut.
\end{itemize}

The example shows:

\begin{itemize}
    \item The output cut is $\{a, c, e\}$.
    \item But the minimal cut is either $\{b, e\}$ or $\{a, d\}$.
\end{itemize}

Thus:

Random selection may or may not preserve the minimum cut.

\subsection{Pseudocode (Exact from Slide)}

\subsubsection*{Algorithm 3: Recursive-Randomized-Min-Cut(G, $\alpha$)}

Input:
\begin{itemize}
    \item An undirected multigraph $G$ with $n$ vertices
    \item An integer constant $\alpha > 0$
\end{itemize}

Output:
\begin{itemize}
    \item A cut $C$ of $G$
\end{itemize}

\begin{verbatim}
if n ≤ α^3 then
    C ← a min-cut of G found using brute force (exhaustive) search;
else
    for i ← 1 to α do
        G' ← multigraph obtained by applying n − ⌈n/√α⌉ random contraction steps on G;
        C' ← Recursive-Randomized-Min-Cut(G', α);
        if i = 1 or |C'| < |C| then
            C ← C';
return C;
\end{verbatim}

\section{Comparison: Deterministic vs Randomized Min-Cut}

The slides state:

\begin{quote}
Any specific problem is responsible for deciding the approach.
\end{quote}

\subsection*{Deterministic Min-Cut}

\begin{itemize}
    \item Always guarantees an exact minimum cut.
    \item May have higher time complexity for large graphs.
\end{itemize}

Stoer-Wagner time complexity:

\[
O(V \cdot E + V^2 \log V)
\]

Where:
\begin{itemize}
    \item $V$ = number of vertices
    \item $E$ = number of edges
\end{itemize}

\subsection*{Randomized Min-Cut}

\begin{itemize}
    \item Provides an approximate minimum cut with high probability.
\end{itemize}

Karger’s algorithm time complexity:

\[
O(V^2)
\]

Where:
\begin{itemize}
    \item $V$ = number of vertices
\end{itemize}

\section{Theorem for Min-Cut Set}

The slide states:

\[
\text{The algorithm outputs a min-cut set with probability at least } \frac{2}{n(n-1)}.
\]

Thus:

Success probability $\ge \frac{2}{n(n-1)}$.

\section{Python Simulation}

The slide instructs:

\begin{itemize}
    \item Open Campuswire post regarding today’s lecture (L10)
    \item Download the \texttt{.ipynb} file pasted in the comments
\end{itemize}

\section*{Conclusion}

This lecture covered:

\begin{enumerate}
    \item Definition and applications of min-cut.
    \item Edge contraction as the main operation.
    \item Successful vs unsuccessful runs.
    \item Max-Flow Min-Cut theorem.
    \item Stoer–Wagner deterministic algorithm and pseudocode.
    \item Karger’s randomized algorithm and recursive pseudocode.
    \item Comparison of deterministic and randomized approaches.
    \item Probability guarantee theorem for min-cut.
\end{enumerate}

This completes the faithful reconstruction of Lecture 10 based strictly on the provided slides.

\end{document}
